% Recuperación expositiva de INTERRELACIONES_ESTRUCTURALES_VACIO.md, §§6--8.
% Insertar después de la inversión cúbica y antes de las coordenadas logarítmicas.
\subsection{Composition transportée des réponses constitutives}
\label{amp-comp-sec}

Le transport angulaire et les incidences $90$ et $120$ proviennent de la
même structure APP--TRIT--TPK : les feuillets conservent les évaluations et
leurs retenues ; le régime tritique conserve l’orientation ; le calendrier
TPK transporte la phase et la mémoire jusqu’aux canaux angulaires. La réponse
constitutive applique à ces canaux la fonction $f$ de
\eqref{iii:eq:monotone}. Son inversion permet d’exprimer, dans les coordonnées
de réponse, la composition des transports du même repère. On conserve
ainsi l’ordre constructif : transport, composition et évaluation spectrale.
L’opération obtenue agit sur cette représentation de l’état enrichi ;
son inverse retrouve les canaux, tandis que l’histoire complète de mémoire
demeure dans la structure de provenance.

Soient $s=q^{30}$ et $\psi=f^{-1}$. L’exposant $30$ est le plus grand commun
diviseur des incidences $90$ et $120$. L’inversion cubique détermine
$\psi:(3/4,1)\longrightarrow(0,1)$. En adjoignant la réponse de frontière
$f(1)=3/4$, nous étendons la bijection à
\[
 \psi:D\longrightarrow(0,1],\qquad D=[3/4,1).
\]
Sur cette frontière, on utilise la fonction rationnelle $f$ ; le quotient
$(1-q^{90})/(1-q^{120})$ y a une singularité éliminable.

\begin{proposition}[Monoïde des réponses et coordonnée additive]
\label{amp-comp-monoid}
L’opération
\begin{equation}
 r\odot t=f\bigl(\psi(r)\psi(t)\bigr),\qquad r,t\in D,
 \label{amp-comp-law}
\end{equation}
fait de $D$ un monoïde commutatif et simplifiable d’élément neutre $3/4$.
L’application
\begin{equation}
 \ell(r)=-\frac1{30}\log\psi(r)
 \label{amp-comp-character}
\end{equation}
est un isomorphisme de ce monoïde avec $(\mathbb R_{\geq0},+)$.
L’intervalle $(3/4,1)$ est stable par $\odot$ ; le neutre appartient à
l’extension de frontière. Le seul élément inversible dans $D$
est le neutre.
\end{proposition}
\begin{proof}
Le produit de deux éléments de $(0,1]$ appartient au même intervalle et
\[
 \psi(r\odot t)=\psi(r)\psi(t).
\]
Pour trois réponses, les deux ordres d’association ont pour image
$\psi(r)\psi(t)\psi(u)$ ; l’injectivité de $\psi$ prouve
l’associativité. La commutativité provient du produit scalaire. Si
$r\odot t=r\odot u$, la positivité de $\psi(r)$ permet de simplifier et
d’obtenir $\psi(t)=\psi(u)$, donc $t=u$. Le paramètre $1$ correspond à
$3/4$ et fournit l’identité. Le logarithme transforme le produit en
somme, et son image dans $(0,1]$ prouve que $\ell$ est une bijection sur
$\mathbb R_{\geq0}$. En particulier, pour la réponse d’un canal $q$,
$\ell(r)=-\log q$. Si les deux paramètres sont strictement inférieurs à
$1$, leur produit l’est aussi. Enfin, deux nombres de $(0,1]$
ont pour produit $1$ uniquement s’ils sont tous deux $1$.
\end{proof}

\begin{proposition}[Simplification admissible et prolongement inverse]
\label{amp-comp-inverse}
Pour $r,u\in D$, l’équation $r\odot t=u$ a une solution $t\in D$
exactement lorsque $u\geq r$. Cette solution est unique et donnée par
\begin{equation}
 t=f\!\left(\frac{\psi(u)}{\psi(r)}\right).
 \label{amp-comp-cancellation}
\end{equation}
La même fonction $f$ est une bijection de $(0,\infty)$ sur $(0,1)$.
Avec cette extension de $\psi$, la loi~\eqref{amp-comp-law} transporte le
groupe multiplicatif des paramètres positifs sur l’intervalle $(0,1)$.
L’inversion du groupe satisfait
\begin{equation}
 r^{\ominus}=\psi(r)r,\qquad
 r\odot r^{\ominus}=\frac34.
 \label{amp-comp-group-inverse}
\end{equation}
\end{proposition}
\begin{proof}
L’équation équivaut à
$\psi(t)=\psi(u)/\psi(r)$. Ce quotient appartient à $(0,1]$
exactement lorsque $\psi(u)\leq\psi(r)$, ce qui, par monotonie décroissante,
équivaut à $u\geq r$. La bijection prouve l’existence et l’unicité.
La dérivée de \eqref{iii:eq:monotone} est négative pour tout $s>0$ ;
les limites de $f$ en $0$ et à l’infini sont $1$ et $0$, respectivement.
Cela établit la bijection étendue. En multipliant le numérateur et le
dénominateur par $s^3$, on obtient
\[
 f(s^{-1})=\frac{s^3+s^2+s}{s^3+s^2+s+1}=s f(s).
\]
L’inverse multiplicatif de $s=\psi(r)$ a donc pour réponse
$\psi(r)r$. Sa composition avec $r$ correspond au paramètre $1$.
\end{proof}

Le prolongement à $s>1$, ou de manière équivalente à $q>1$, est l’extension
algébrique du lecteur. Le secteur physique contractant du chapitre conserve
son domaine déclaré. La coordonnée $\ell$ se prolonge en un isomorphisme
avec $(\mathbb R,+)$ et change de signe sous l’inversion du groupe.

\subsubsection{Réalisation opératorielle dans un repère de feuillets commun}

\begin{proposition}[Composition avec projecteurs marqués]
\label{amp-comp-operator}
Soient $\mathcal R=\mathcal R^*$, $\mathcal R^2=I$ et
$P_\pm=(I\pm\mathcal R)/2$. Pour $j=a,b$, soient
\[
 T_j=q_{j,+}P_++q_{j,-}P_-,\qquad
 \mathcal V_j=f(T_j^{30}),\qquad 0<q_{j,\pm}<1.
\]
La composition $T_{a\circ b}=T_aT_b$ satisfait
\begin{equation}
 \mathcal V_{a\circ b}
 =f\bigl(\psi(\mathcal V_a)\psi(\mathcal V_b)\bigr).
 \label{amp-comp-operator-law}
\end{equation}
Ses deux réponses sont $r_{a,\pm}\odot r_{b,\pm}$. Si
$T_j=\exp[-(x_jI+y_j\mathcal R)]$ avec $x_j>|y_j|$, alors
\begin{equation}
 T_aT_b=
 \exp[-((x_a+x_b)I+(y_a+y_b)\mathcal R)],
 \label{amp-comp-angular-addition}
\end{equation}
et $x_a+x_b>|y_a+y_b|$.
\end{proposition}
\begin{proof}
L’orthogonalité $P_+P_-=0$ donne
\[
 T_aT_b=q_{a,+}q_{b,+}P_++q_{a,-}q_{b,-}P_-.
\]
En particulier, $T_a$ et $T_b$ commutent. Le calcul fonctionnel dans ces
projecteurs fournit
$\psi(\mathcal V_j)=T_j^{30}$ et
$(T_aT_b)^{30}=T_a^{30}T_b^{30}$ ; appliquer $f$ prouve
\eqref{amp-comp-operator-law}. Pour toute fonction $g$ définie sur les
deux valeurs propres, on a, plus explicitement,
\[
 P_\pm g(T_j)=g(q_{j,\pm})P_\pm.
\]
Comme $q_{j,\pm}=\exp[-(x_j\pm y_j)]$, leurs produits prouvent
\eqref{amp-comp-angular-addition}. L’inégalité triangulaire conclut
$x_a+x_b>|y_a|+|y_b|\geq|y_a+y_b|$.
\end{proof}

Les coordonnées retrouvées peuvent s’écrire directement comme
\[
 x=\frac{\ell(r_+)+\ell(r_-)}2,\qquad
 y=\frac{\ell(r_+)-\ell(r_-)}2.
\]
La composition additionne ces paires. L’échange simultané de feuillets
permute les deux canaux de chaque transport, agit comme
$(x,y)\mapsto(x,-y)$ et est un automorphisme de la loi composante par
composante. L’inversion des deux transports dans l’extension algébrique
agit comme $(x,y)\mapsto(-x,-y)$. Les deux involutions commutent et
conservent des fonctions distinctes : l’une échange les marques d’orientation ;
l’autre inverse le transport.

Pour représenter l’échange par conjugaison, on utilise un
opérateur $L$ satisfaisant $L\mathcal R L^{-1}=-\mathcal R$. Dans la
représentation
\[
 \mathcal R=\begin{pmatrix}0&1\\1&0\end{pmatrix},\qquad
 L=\begin{pmatrix}1&0\\0&-1\end{pmatrix},
\]
la multiplication directe donne $LP_\pm L^{-1}=P_\mp$ et, par conséquent,
$LT_jL^{-1}=q_{j,-}P_++q_{j,+}P_-$. La conjugaison par
$\mathcal R$ conserve ses propres projecteurs. Une trace qui réunit les
canaux ne remplace pas non plus les projecteurs dans cette identité fonctionnelle :
l’information de leurs étiquettes intervient dans la composition.

L’hypothèse de repère commun est substantielle. Par exemple, les opérateurs
positifs définis
\[
 A_0=\begin{pmatrix}1/2&0\\0&1/3\end{pmatrix},\qquad
 B_0=\begin{pmatrix}1/2&1/6\\1/6&1/2\end{pmatrix}
\]
ont des valeurs propres positives, mais
\[
 A_0B_0=\begin{pmatrix}1/4&1/12\\1/18&1/6\end{pmatrix}
\]
n’est pas autoadjoint. La proposition spécifie la composition des
transports dans les projecteurs communs ; son domaine ne comprend pas un mélange
arbitraire de milieux matériels. La réalisation d’une concaténation de
chemins HMT conserve en outre ses données de mémoire : la fermeture algébrique du
cône de paramètres décrit la représentation et ne sélectionne pas à elle seule
de nouveaux chemins primaires.

\subsubsection{Produit de canaux et composition de transports}

La vitesse normalisée d’un état reste déterminée par
\[
 \widehat c=(r_+r_-)^{-1}.
\]
Ici, le produit ordinaire réunit les deux réponses de ce même état.
La loi $\odot$ évalue une autre opération : la composition de deux transports
dans chaque feuillet. Le calcul exact suivant distingue leurs arguments :
\begin{equation}
 f(1/2)=\frac{14}{15},\qquad
 \frac{14}{15}\odot\frac{14}{15}=f(1/4)=\frac{84}{85}
 \ne\frac{196}{225}=\left(\frac{14}{15}\right)^2.
 \label{amp-comp-rational-example}
\end{equation}
Les fractions de cet exemple sont des paramètres algébriques de vérification.
L’expression de $\widehat c$ et les identités constitutives précédentes
restent inchangées. Le contenu ajouté est une loi explicite pour
transporter la composition et retrouver sa coordonnée angulaire additive.
